% ================================================================================
% RUCP SUBMISSION VERSION — ~10,000 palabras
% Revista Uruguaya de Ciencia Política
% ================================================================================
% Basado en: Informe_Electoral_Uruguay_2019_2024
% Normas RUCP: https://www.scielo.edu.uy/scielo.php?script=sci_serial&pid=1688-499X
% Máximo: 10,000 palabras (texto principal, sin referencias ni tablas)
% Figuras: máximo 8 en texto principal
% ================================================================================

\documentclass[12pt,a4paper]{article}

% ── Codificación y idioma ────────────────────────────────────────────────────
\usepackage[utf8]{inputenc}
\usepackage[spanish]{babel}
\usepackage[margin=2.5cm]{geometry}

% ── Tipografía ───────────────────────────────────────────────────────────────
\usepackage{lmodern}
\usepackage{microtype}
\usepackage[onehalfspacing]{setspace}

% ── Matemáticas ─────────────────────────────────────────────────────────────
\usepackage{amsmath}
\usepackage{amssymb}

% ── Gráficos y figuras ───────────────────────────────────────────────────────
\usepackage{graphicx}
\usepackage{float}
\usepackage[labelfont=bf,font=small]{caption}

% ── Tablas ───────────────────────────────────────────────────────────────────
\usepackage{booktabs}
\usepackage{multirow}
\usepackage{adjustbox}
\usepackage{siunitx}
\renewcommand{\arraystretch}{1.1}

% ── Control de floats ────────────────────────────────────────────────────────
\usepackage[section]{placeins}
\renewcommand{\topfraction}{0.85}
\renewcommand{\textfraction}{0.15}
\renewcommand{\floatpagefraction}{0.7}

% ── Párrafos ─────────────────────────────────────────────────────────────────
\setlength{\parindent}{15pt}
\setlength{\parskip}{0pt plus 1pt}
\widowpenalty=10000
\clubpenalty=10000

% ── Citas y referencias ──────────────────────────────────────────────────────
\usepackage{natbib}
\bibliographystyle{apalike}

% ── Hipervínculos ────────────────────────────────────────────────────────────
\usepackage{hyperref}
\hypersetup{
    colorlinks=true,
    linkcolor=black,
    citecolor=black,
    urlcolor=black,
    pdftitle={Rebelión rural en el balotaje: heterogeneidad geográfica en transferencias de votos, Uruguay 2019-2024},
    pdfauthor={[Autor/a]},
}

% ── Rutas ────────────────────────────────────────────────────────────────────
\graphicspath{{../../outputs/figures/}}
\newcommand{\tablepath}{../../outputs/tables/latex/}

% ── Comandos personalizados ──────────────────────────────────────────────────
\newcommand{\rhat}{$\hat{R}$}
\newcommand{\CI}[2]{[#1, #2]}

% ================================================================================
% METADATOS
% ================================================================================
\title{%
    \textbf{Rebelión rural en el balotaje: heterogeneidad geográfica}\\
    \textbf{en transferencias de votos, Uruguay 2019--2024}\\[0.5em]
    {\normalsize \textit{Versión para evaluación — RUCP}}
}
\author{[Autor/a]\thanks{[Institución. Correspondencia: email@inst.edu.uy]}}
\date{}

% ================================================================================
\begin{document}
% ================================================================================

\maketitle

% ── ABSTRACT ────────────────────────────────────────────────────────────────
\begin{abstract}
\noindent
Se documenta un patrón contraintuitivo en el balotaje presidencial uruguayo de 2024: los circuitos rurales exhibieron tasas de defección de Cabildo Abierto (CA) al Frente Amplio (FA) 21{,}8 puntos porcentuales más altas que los circuitos urbanos (60{,}8\% vs.\ 39{,}0\%), contradiciendo la expectativa de mayor conservadurismo rural para un partido de perfil militar-populista. Este hallazgo constituye la tesis central del estudio. Utilizando un modelo bayesiano de transferencia de votos con prior Dirichlet sobre 7{,}271 circuitos electorales en 2024 y 7{,}118 en 2019, se estiman matrices de transición nacionales, departamentales y estratificadas. A nivel nacional, el 46{,}5\% [IC\,95\%: 43{,}0\%, 49{,}9\%] de los votantes de CA transfirió su voto al FA en el balotaje. La brecha geográfica es estadísticamente robusta: $P(T_{\text{rural}} > T_{\text{urbano}} \mid \text{datos}) > 0{,}99$. La comparación temporal 2019--2024 revela que, pese a mejoras generalizadas en la cohesión coalicional, la coalición perdió por el colapso del caudal de CA ($-$78\% entre elecciones). El análisis ofrece evidencia de una reconfiguración geográfica profunda del electorado uruguayo no capturada por análisis nacionales agregados.
\end{abstract}

\vspace{0.4cm}
\noindent\textbf{Palabras clave:} heterogeneidad geográfica, brecha urbano-rural, transferencias de votos, inferencia ecológica bayesiana, elecciones Uruguay 2024, Cabildo Abierto, balotaje.

\vspace{0.3cm}
\noindent\textbf{Keywords:} geographic heterogeneity, urban-rural divide, vote transfers, Bayesian ecological inference, Uruguay 2024 elections, second-round ballot.

\newpage
\tableofcontents
\newpage

% ================================================================================
% SECCIONES PRINCIPALES
% ================================================================================

% =========================================================
% INTRODUCCIÓN — versión RUCP (~1,500 palabras)
% =========================================================

\section{Introducción}
\label{sec:introduccion}

El 24 de noviembre de 2024, Yamandú Orsi (Frente Amplio, FA) derrotó a Álvaro Delgado (Partido Nacional, PN) en el balotaje presidencial uruguayo por aproximadamente 95{,}000 votos, poniendo fin a cinco años de gobierno de la Coalición Republicana. El resultado sorprendió a analistas que estimaban que la coalición ---integrada por el PN, el Partido Colorado (PC), Cabildo Abierto (CA) y el Partido Independiente (PI)--- capitalizaría su organización en segunda vuelta. Un análisis posterior revela un patrón inesperado que motiva este estudio: los votantes rurales de CA transfirieron su apoyo al FA en proporciones \textit{mayores} que los votantes urbanos del mismo partido, contradiciendo la expectativa convencional de mayor conservadurismo en zonas rurales para un partido de perfil militar-populista. Se denomina este fenómeno ``rebelión rural'' y constituye la tesis central del trabajo.

La evidencia de este patrón no es accesible de forma directa. El secreto del voto impide conocer las decisiones individuales; solo se dispone de recuentos de votos agregados a nivel de circuito electoral. El desafío de inferir comportamiento individual a partir de datos agregados ---el \textit{problema de la inferencia ecológica} \citep{king1997}--- es la razón por la que la \textit{falacia ecológica} \citep{robinson1950ecological} acecha a los análisis de transferencias que utilizan correlaciones agregadas sin corrección estadística. El trabajo fundacional de \citet{king1997} propuso un marco bayesiano para resolver este problema en el caso binario; \citet{rosen2001bayesian} extendieron el enfoque al caso multidimensional $R \times C$ requerido por sistemas multipartidistas como el uruguayo. Trabajo reciente de \citet{greiner2009rxc} evaluó los supuestos de identificación de estos modelos, e \citet{imai2008bayesian} propusieron formulaciones alternativas. En América Latina, las aplicaciones sistemáticas de estos métodos a nivel circuital siguen siendo escasas \citep{calvo2003local}, confiriéndole al presente estudio un carácter exploratorio regionalmente relevante.

El caso de CA merece atención especial en la literatura. Fundado en 2019 por el general retirado Guido Manini Ríos, irrumpió en el sistema de partidos con el 11\% de los votos y colapsó al 2.7\% en 2024 ($-$78\% en caudal electoral). La comprensión del destino de sus votantes remanentes en el balotaje ---y de por qué los circuitos rurales exhiben patrones distintos--- resulta decisiva para interpretar la estabilidad del sistema de partidos uruguayo, uno de los más institucionalizados de América Latina \citep{mainwaring1995party, sartori1976parties}. La consolidación del FA como actor central del sistema ha sido documentada por \citet{perezbentancur2019frente}; el rol de las alianzas coalicionales en el presidencialismo latinoamericano ha sido analizado por \citet{chasquetti2004democracia}.

La literatura sobre divisiones urbano-rurales sugiere que los votantes rurales tienden a preferir partidos de derecha, particularmente en contextos donde la economía agrícola es central \citep{rodden2019s, gonzalez1991electoral}. La segunda vuelta, sin embargo, introduce dinámicas propias: la elección binaria fuerza un reposicionamiento estratégico que puede revertir alineaciones habituales \citep{birch2003tworound}. Cuando el partido del votante queda fuera del balotaje ---como le ocurrió a todos los socios menores de la coalición--- la decisión de cuál de los dos candidatos apoyar queda determinada por factores sustantivos (política pública, liderazgos, identidades) que pueden diferir sistemáticamente entre contextos geográficos.

El presente estudio realiza las siguientes contribuciones:

\begin{enumerate}
    \item \textbf{Documenta la ``rebelión rural''}: los circuitos rurales exhiben tasas de defección de CA al FA 21{,}8 pp superiores a los urbanos ($P(T_{\text{rural}} > T_{\text{urbano}} \mid \text{datos}) > 0{,}99$), el hallazgo más contraintuitivo del estudio.

    \item \textbf{Primera aplicación sistemática de inferencia ecológica bayesiana a nivel circuital en balotajes uruguayos}: 7{,}271 circuitos en 2024 y 7{,}118 en 2019, máxima resolución geográfica disponible en datos oficiales.

    \item \textbf{Identifica la paradoja de la coalición cohesiva derrotada}: pese a mejoras generalizadas en la retención de votos entre 2019 y 2024, la coalición perdió el balotaje ---efecto directo del colapso del caudal de CA ($-$78\%)---, con implicancias para la teoría de coaliciones presidenciales.

    \item \textbf{Valida metodológicamente los resultados}: análisis de sensibilidad con priors Dirichlet genuinamente distintos ($\alpha \in \{0.5, 1, 5\}$) y 500 replicaciones bootstrap confirman la estabilidad de los hallazgos.
\end{enumerate}

El artículo se organiza como sigue. La Sección~\ref{sec:metodologia} describe el modelo bayesiano de transferencia de votos, los datos y la estrategia analítica. La Sección~\ref{sec:resultados} presenta los resultados nacionales y la heterogeneidad geográfica. La Sección~\ref{sec:temporal} examina la evolución 2019--2024. La Sección~\ref{sec:discusion} discute las implicaciones teóricas y las limitaciones del estudio.

\newpage

% =========================================================
% METODOLOGÍA — versión RUCP (~1,500 palabras)
% =========================================================

\section{Metodología}
\label{sec:metodologia}

\subsection{Datos}

Se utilizan datos electorales oficiales de la Corte Electoral de Uruguay para dos pares de elecciones: primera vuelta y balotaje de 2024 (27 de octubre y 24 de noviembre) y de 2019 (27 de octubre y 24 de noviembre). La unidad de análisis es el \textit{circuito electoral}, la unidad geográfica de menor desagregación en los datos oficiales. El conjunto de datos 2024 abarca 7{,}271 circuitos; el de 2019, 7{,}118.

Para la primera vuelta se clasifican los votos en seis categorías: FA, PN, PC, CA, PI y OTROS (partidos menores agrupados). Para el balotaje, las opciones son FA, PN (representando a la coalición) y votos en blanco/nulos. Se aplica un umbral mínimo de 30 votos válidos por elección para excluir circuitos con datos insuficientes. Los controles de calidad verifican que los totales a nivel de circuito coincidan con los agregados departamentales oficiales; menos del 0.1\% de los circuitos presentó valores faltantes (eliminados por lista).

La estratificación geográfica distingue circuitos urbanos (86.8\%, $n=6{,}312$) de rurales (13.2\%, $n=959$), y el Área Metropolitana (Montevideo y Canelones, $\approx$60\% de la población) del Interior (17 departamentos restantes). La clasificación urbano/rural proviene de datos administrativos electorales.

\subsection{El Modelo Bayesiano de Transferencia de Votos}

Se implementa un modelo bayesiano de transferencia de votos con prior Dirichlet (complete pooling), inspirado en el marco de inferencia ecológica de \citet{king1997} y su extensión multidimensional por \citet{rosen2001bayesian}. El modelo estima una matriz de transición $\mathbf{T}$ de dimensión $K \times J$ ($K$ partidos de origen, $J$ de destino), donde $T_{kj}$ representa la probabilidad de que un votante que apoyó al partido $k$ en la primera vuelta haya votado por la opción $j$ en el balotaje.

Para cada circuito $i$, sea $\mathbf{X}_i$ el vector de votos de primera vuelta e $\mathbf{Y}_i$ el de votos en el balotaje. El modelo especifica:
\begin{equation}
  \mathbf{T}_k \sim \text{Dirichlet}(\boldsymbol{\alpha}), \quad \alpha_j = 1\ \forall j \qquad \text{(prior uniforme)}
  \label{eq:prior}
\end{equation}
\begin{equation}
  \mathbf{Y}_i \sim \mathcal{N}\!\left(\mathbf{X}_i \mathbf{T},\; \sigma^2\right)
  \label{eq:likelihood}
\end{equation}
donde $\sigma_{ij} = \sqrt{\mu_{ij} + 1}$ es la desviación estándar de la aproximación Normal a la distribución Multinomial \citep{martin2021bayesian}. El prior Dirichlet garantiza $T_{kj} \in [0,1]$ y $\sum_j T_{kj} = 1$ para cada partido de origen.

\textbf{Nota sobre nomenclatura}: este modelo difiere del modelo de King (1997) en un aspecto fundamental. King propuso un modelo con parámetros que varían por unidad geográfica (partial pooling). El presente modelo estima una única $\mathbf{T}$ para todos los circuitos (complete pooling), lo que es computacionalmente eficiente con 7{,}271 observaciones y produce estimaciones nacionales robustas. El modelo jerárquico de partial pooling queda propuesto como extensión futura. Adicionalmente, se estima un segundo modelo con verosimilitud DirichletMultinomial \citep{mosimann1962compound}, estadísticamente más apropiado para datos de conteo, que confirma la dirección e interpretación de los hallazgos principales.

\subsection{Estimación y Diagnósticos}

La inferencia posterior se realiza mediante MCMC con el No-U-Turn Sampler (NUTS) \citep{betancourt2017conceptual} implementado en PyMC \citep{abril2023pymc}, utilizando el backend \texttt{nutpie} (Rust/numba) para acelerar la exploración del espacio de parámetros. El análisis nacional utiliza 4 cadenas $\times$ 4{,}000 muestras post-calentamiento (2{,}000 warmup, tasa de aceptación objetivo 0.95); los análisis estratificados emplean 4 cadenas $\times$ 2{,}000 muestras.

La convergencia se evalúa mediante el estadístico de Gelman-Rubin ($\hat{R}$) y el tamaño de muestra efectivo (ESS) \citep{vehtari2021rhat, gelman1992inference}. Se exigen $\hat{R} < 1.01$ y $\text{ESS} > 1{,}000$ para todos los parámetros; solo los modelos que satisfacen ambos criterios se incluyen en los resultados. Todos los modelos reportados convergieron exitosamente.

\subsection{Validación y Análisis de Sensibilidad}

Se realizaron tres tipos de validación. Primero, verificaciones predictivas posteriores (PPC) comparan la distribución de votos simulados con la observada. Segundo, se verificó que las predicciones a nivel circuital reproducen los agregados departamentales y nacionales al ponderar por totales de votos. Tercero, se implementó un análisis de sensibilidad que compara tres especificaciones de prior Dirichlet ($\alpha \in \{0.5, 1.0, 5.0\}$) y 500 replicaciones bootstrap (remuestreo del 80\% de circuitos con reemplazo). La diferencia máxima entre especificaciones de prior es de 1.3 pp; el coeficiente de variación bootstrap para CA$\to$FA es del 6.4\% (error estándar del bootstrap $< 0.2$ pp).

\subsection{Comparación Geográfica Formal}

Para probar formalmente si las tasas de defección difieren entre estratos geográficos, se computan diferencias de posteriors muestra a muestra:
\begin{equation}
P\!\left(T_{kj}^{\text{rural}} > T_{kj}^{\text{urbano}} \mid \text{datos}\right) = \frac{1}{S}\sum_{s=1}^S \mathbf{1}\!\left[T_{kj,s}^{\text{rural}} > T_{kj,s}^{\text{urbano}}\right]
\end{equation}
Este enfoque aprovecha directamente el marco bayesiano y produce evidencia cuantificada sobre la dirección de las diferencias geográficas, sin depender de la superposición de intervalos de credibilidad como proxy \citep{gelman2013bayesian}.

\newpage

% =========================================================
% RESULTADOS — versión RUCP (~2,500 palabras)
% Combina: nacionales + urbano/rural + metro/interior + departamental (mapa)
% =========================================================

\section{Resultados}
\label{sec:resultados}

\subsection{Estimaciones Nacionales}
\label{sec:resultados_nacional}

El modelo ajustado sobre los 7{,}271 circuitos electorales de 2024 convergió exitosamente ($\hat{R}_{\max} = 1.001$, $\text{ESS}_{\min} > 6{,}000$). La Tabla~\ref{tab:national_matrix_pi} presenta la matriz de transición nacional completa con intervalos de credibilidad del 95\%.

\begin{table}[htbp]
\centering
\caption{Matriz de Transición Nacional (Primera Vuelta $\rightarrow$ Ballotage) con PI}
\label{tab:national_matrix_pi}
\small
\begin{tabular}{lS[table-format=2.1]S[table-format=2.1]S[table-format=2.1]}
\toprule
{\textbf{Origen}} & {\textbf{FA (\%)}} & {\textbf{PN (\%)}} & {\textbf{Blancos (\%)}} \\
\midrule
CA & 46.5 & 49.1 & 4.4 \\
 & {[43.0, 49.9]} & {[45.6, 52.7]} & {[3.4, 5.5]} \\
FA & 98.9 & 0.5 & 0.6 \\
 & {[98.6, 99.2]} & {[0.2, 0.8]} & {[0.5, 0.7]} \\
OTROS & 21.9 & 69.4 & 8.7 \\
 & {[18.4, 25.4]} & {[65.9, 72.8]} & {[7.8, 9.7]} \\
PC & 10.0 & 87.9 & 2.1 \\
 & {[9.1, 11.0]} & {[86.9, 88.8]} & {[1.8, 2.4]} \\
PI & 0.0 & 99.9 & 0.1 \\
 & {[0.0, 0.2]} & {[99.7, 100.0]} & {[0.0, 0.2]} \\
PN & 6.3 & 91.8 & 2.0 \\
 & {[5.8, 6.8]} & {[91.2, 92.3]} & {[1.8, 2.1]} \\
\bottomrule
\end{tabular}
\begin{flushleft}
\small \textit{Nota:} CA: Cabildo Abierto; FA: Frente Amplio; PC: Partido Colorado; PI: Partido Independiente; PN: Partido Nacional. Los valores representan medias posteriores con intervalos creíbles al 95\% entre corchetes. Max $\hat{R} < 1.01$ indica convergencia MCMC adecuada.
\end{flushleft}
\end{table}


\textbf{Cabildo Abierto} exhibió la defección más pronunciada: se estima que el 46.5\% [IC\,95\%: 43.0\%, 49.9\%] de sus votantes de primera vuelta transfirió el voto al FA en el balotaje, mientras que el 49.1\% [45.6\%, 52.7\%] apoyó al PN. Esta división casi igualitaria contrasta con la expectativa de que un partido de perfil conservador-nacional retuviera ampliamente su voto en favor de la coalición.

\textbf{Partido Colorado}: 87.9\% [86.9\%, 88.8\%] fue al PN; el 10.0\% [9.1\%, 11.0\%] al FA. En términos absolutos, las 387{,}125 bases del PC implican $\approx$38{,}700 votos perdidos hacia FA.

\textbf{Partido Nacional}: 91.8\% [91.2\%, 92.3\%] se autoconservó; 6.3\% [5.8\%, 6.8\%] fue al FA ($\approx$40{,}700 votos). \textbf{Partido Independiente}: retención casi perfecta en el PN (99.9\%), defección mínima al FA ($\approx$0\%).

En conjunto, la coalición transfirió al FA un estimado de $\approx$107{,}000 votos, superando el margen de derrota de $\approx$95{,}000 votos. La lealtad del FA fue asimétrica: 98.9\% [98.6\%, 99.2\%] de sus votantes de primera vuelta permaneció en el FA. La conclusión es que el resultado del balotaje no se explica por una movilización del FA, sino por las fracturas internas de la coalición.

\subsection{La Rebelión Rural: Heterogeneidad Urbano-Rural}
\label{sec:resultados_rural}

El hallazgo central del estudio emerge de la estratificación por clasificación urbano/rural. La Tabla~\ref{tab:urban_rural_pi} presenta las estimaciones para ambos estratos.

\begin{table}[htbp]
\centering
\caption{Tasas de Defección de Partidos de la Coalición por Clasificación Urbana/Rural (con PI)}
\label{tab:urban_rural_pi}
\small
\begin{tabular}{lS[table-format=4]S[table-format=2.1]S[table-format=2.1]S[table-format=2.1]S[table-format=2.1]S[table-format=2.1]S[table-format=2.1]S[table-format=2.1]S[table-format=2.1]}
\toprule
{\textbf{Estrato}} & {\textbf{N}} & \multicolumn{2}{c}{\textbf{CA}} & \multicolumn{2}{c}{\textbf{PC}} & \multicolumn{2}{c}{\textbf{PI}} & \multicolumn{2}{c}{\textbf{PN}} \\
\cmidrule(lr){3-4} \cmidrule(lr){5-6} \cmidrule(lr){7-8} \cmidrule(lr){9-10}
 & & {\textbf{$\rightarrow$FA (\%)}} & {\textbf{$\rightarrow$PN (\%)}} & {\textbf{$\rightarrow$FA (\%)}} & {\textbf{$\rightarrow$PN (\%)}} & {\textbf{$\rightarrow$FA (\%)}} & {\textbf{$\rightarrow$PN (\%)}} & {\textbf{$\rightarrow$FA (\%)}} & {\textbf{$\rightarrow$PN (\%)}} \\
\midrule
Urbano & 6311 & 39.0 & 55.5 & 5.2 & 92.5 & 0.1 & 99.9 & 6.9 & 91.2 \\
 & & {[3510.0, 4280.0]} & {[5170.0, 5940.0]} & {[410.0, 630.0]} & {[9200.0, 9300.0]} & {[0.0, 20.0]} & {[9970.0, 10000.0]} & {[630.0, 750.0]} & {[9010.0, 9260.0]} \\
Rural & 960 & 60.8 & 38.2 & 18.2 & 80.3 & 0.8 & 98.2 & 7.0 & 91.1 \\
 & & {[5170.0, 6980.0]} & {[2910.0, 4730.0]} & {[1620.0, 2010.0]} & {[8600.0, 8820.0]} & {[0.0, 50.0]} & {[9920.0, 10000.0]} & {[580.0, 820.0]} & {[8900.0, 9370.0]} \\
\bottomrule
\end{tabular}
\begin{flushleft}
\small \textit{Nota:} CA: Cabildo Abierto; PC: Partido Colorado; PI: Partido Independiente; PN: Partido Nacional; FA: Frente Amplio. Los valores representan medias posteriores con intervalos creíbles al 95\% entre corchetes. Max $\hat{R} < 1.01$ indica convergencia MCMC adecuada.
\end{flushleft}
\end{table}


\textbf{Cabildo Abierto}: los circuitos rurales exhiben una defección CA$\to$FA del 60.8\% [51.7\%, 69.8\%], frente al 39.0\% [35.1\%, 42.8\%] en los urbanos. La brecha de 21.8 puntos porcentuales es estadísticamente robusta: $P(T_{\text{CA}\to\text{FA}}^{\text{rural}} > T_{\text{CA}\to\text{FA}}^{\text{urbano}} \mid \text{datos}) > 0{,}99$. Los intervalos de credibilidad del 95\% no se superponen, confirmando que la diferencia no es atribuible a incertidumbre de muestreo MCMC.

\textbf{Partido Colorado}: los circuitos rurales muestran una defección PC$\to$FA del 18.2\% [7.1\%, 29.2\%] frente al 7.2\% [5.7\%, 8.7\%] urbano. La brecha (11.0 pp) es también estadísticamente significativa pero con mayor incertidumbre en el estrato rural, reflejo del menor número de circuitos rurales con presencia colorada relevante.

\textbf{Partido Nacional}: retención casi idéntica en ambos estratos (91.6\% rural vs.\ 91.8\% urbano), indicando que los votantes del PN mantuvieron lealtad uniforme independientemente del contexto geográfico.

Este patrón contradice directamente la narrativa convencional de conservadurismo rural en Uruguay \citep{gonzalez1991electoral, rodden2019s}. Una posible explicación es que los circuitos rurales concentran un perfil de votante de CA más volátil: electores que apoyaron al partido por su posicionamiento anti-establishment en 2019, pero cuya desafección con el gobierno de la coalición en 2020--2024 fue más intensa en zonas donde la crisis agrícola (sequía, precios de commodities, costos de producción) impactó de forma directa. El diseño de inferencia ecológica no permite confirmar mecanismos causales, pero la geografía del fenómeno es consistente con esta hipótesis.

\begin{figure}[H]
\centering
\includegraphics[width=0.88\textwidth]{urban_rural_ca_defection_forest.png}
\caption{Comparación de tasas de defección CA$\to$FA por estrato geográfico. Los puntos representan medias posteriores; las barras, intervalos de credibilidad del 95\%. La brecha urbano-rural de 21.8 pp es estadísticamente robusta ($P > 0.99$).}
\label{fig:urban_rural_forest}
\end{figure}

\subsection{Dimensión Metropolitana-Interior}
\label{sec:resultados_metro}

El contraste Área Metropolitana (Montevideo + Canelones) vs.\ Interior refuerza el patrón de mayor volatilidad fuera de la capital.

\begin{table}[htbp]
\centering
\caption{Variación Regional en las Defecciones de la Coalición (con PI)}
\label{tab:region_pi}
\small
\begin{tabular}{lS[table-format=4]S[table-format=2.1]S[table-format=2.1]S[table-format=2.1]S[table-format=2.1]S[table-format=2.1]S[table-format=2.1]S[table-format=2.1]S[table-format=2.1]}
\toprule
{\textbf{Región}} & {\textbf{N}} & \multicolumn{2}{c}{\textbf{CA}} & \multicolumn{2}{c}{\textbf{PC}} & \multicolumn{2}{c}{\textbf{PI}} & \multicolumn{2}{c}{\textbf{PN}} \\
\cmidrule(lr){3-4} \cmidrule(lr){5-6} \cmidrule(lr){7-8} \cmidrule(lr){9-10}
 & & {\textbf{$\rightarrow$FA (\%)}} & {\textbf{$\rightarrow$PN (\%)}} & {\textbf{$\rightarrow$FA (\%)}} & {\textbf{$\rightarrow$PN (\%)}} & {\textbf{$\rightarrow$FA (\%)}} & {\textbf{$\rightarrow$PN (\%)}} & {\textbf{$\rightarrow$FA (\%)}} & {\textbf{$\rightarrow$PN (\%)}} \\
\midrule
Área Metropolitana & 3671 & 23.2 & 69.9 & 3.8 & 94.8 & 0.2 & 99.6 & 0.3 & 97.7 \\
 & & {[17.4, 28.8]} & {--} & {[2.1, 5.2]} & {--} & {[0.0, 0.8]} & {--} & {[0.0, 1.1]} & {--} \\
Interior & 3600 & 54.7 & 42.0 & 10.3 & 87.5 & 0.2 & 99.4 & 6.9 & 91.3 \\
 & & {[50.3, 59.1]} & {--} & {[9.2, 11.3]} & {--} & {[0.0, 0.8]} & {--} & {[6.3, 7.5]} & {--} \\
\bottomrule
\end{tabular}
\begin{flushleft}
\small \textit{Nota:} CA: Cabildo Abierto; PC: Partido Colorado; PI: Partido Independiente; PN: Partido Nacional; FA: Frente Amplio. Área Metropolitana incluye Montevideo y Canelones. Los valores representan medias posteriores con intervalos creíbles al 95\% entre corchetes para transferencias a FA.
\end{flushleft}
\end{table}


La defección CA$\to$FA en el Interior (51.0\% [45.3\%, 56.7\%]) supera en 10.2 pp a la del Área Metropolitana (40.8\% [35.8\%, 45.8\%]). La prueba formal de diferencia posterior confirma $P(T_{\text{CA}\to\text{FA}}^{\text{interior}} > T_{\text{CA}\to\text{FA}}^{\text{metro}} \mid \text{datos}) = 0.96$, evidencia positiva aunque no concluyente al nivel del contraste urbano-rural. El Interior concentra la mayor parte de los circuitos rurales, por lo que ambas estratificaciones capturan dimensiones parcialmente superpuestas de la misma heterogeneidad geográfica.

\subsection{Variación Departamental}
\label{sec:resultados_dept}

La variación entre los 19 departamentos excede toda la heterogeneidad documentada en los estratos anteriores. La Figura~\ref{fig:mapa_dept} muestra el mapa de tasas de defección CA$\to$FA por departamento.

\begin{figure}[H]
\centering
\includegraphics[width=0.92\textwidth]{department_ca_defection_map.png}
\caption{Mapa de tasas de defección CA$\to$FA por departamento (medias posteriores). El rango varía de 6.0\% (Tacuarembó) a 95.6\% (Río Negro), superando en 89.6 pp la heterogeneidad de cualquier estratificación binaria urbano/rural o metropolitana/interior.}
\label{fig:mapa_dept}
\end{figure}

La dispersión departamental ---de 6.0\% (Tacuarembó) a 95.6\% (Río Negro)--- sugiere que la dinámica de la segunda vuelta respondió a factores locales (candidaturas, redes organizacionales, economías regionales específicas) que el modelo nacional de complete pooling promedia pero no captura. Este resultado refuerza la necesidad de modelos jerárquicos como extensión futura.

Las estimaciones para los diez departamentos con mayor y menor variación respecto a la media nacional se presentan en el material suplementario (Tabla~A.1). Las correlaciones interdepartamentales entre 2019 y 2024 son menores a $r < 0.13$, indicando nula estabilidad temporal en la geografía de las defecciones ---una reconfiguración profunda del electorado que no puede explicarse por identidades geográficas persistentes.

\newpage

% =========================================================
% ANÁLISIS TEMPORAL 2019–2024 — versión RUCP (~800 palabras)
% =========================================================

\section{Evolución Temporal: 2019--2024}
\label{sec:temporal}

\subsection{La Paradoja de la Coalición Cohesiva Derrotada}

La comparación entre las elecciones de 2019 y 2024 revela una paradoja: la Coalición Republicana perdió el balotaje de 2024 a pesar de mejorar sustancialmente su cohesión interna respecto a 2019. La Tabla~\ref{tab:comparison_2019_2024} presenta las estimaciones mediante promedios departamentales ponderados por votos, que permiten comparabilidad entre años con cambios en la distribución de circuitos.

\begin{table}[H]
\centering
\caption{Evolución de defecciones coalicionales 2019-2024. Mejora dramática en cohesión: todos los partidos redujeron defección entre 26.6 y 92.1 puntos porcentuales.}
\label{tab:comparison_2019_2024}
\small
\begin{tabular}{lrrrrr}
\toprule
\textbf{Partido} & \textbf{Def. 2019} & \textbf{Def. 2024} & \textbf{Cambio} & \textbf{Votos def.} & \textbf{Votos def.} \\
 & \textbf{(\%)} & \textbf{(\%)} & \textbf{(pp)} & \textbf{2019} & \textbf{2024} \\
\midrule
Cabildo Abierto & 69.5\% & 27.6\% & -41.9 & 186.316 & 16.382 \\
Partido Colorado & 33.8\% & 7.2\% & -26.6 & 101.412 & 27.793 \\
Partido Nacional & 36.0\% & 3.6\% & -32.4 & 250.260 & 23.077 \\
Partido Independiente & 94.8\% & 2.8\% & -92.1 & 22.335 & 1.130 \\
\bottomrule
\multicolumn{6}{l}{\footnotesize Defección = \% de votos transferidos de partido coalicional a FA en balotaje.} \\
\multicolumn{6}{l}{\footnotesize Partido Independiente tuvo el cambio más dramático: -92.1 pp (de 94.8\% a 2.8\%).} \\
\end{tabular}
\end{table}


\textbf{Cabildo Abierto}: en 2019 el 88.7\% de los votantes de CA transfirió su voto al PN en el balotaje; en 2024, esta proporción cayó al 46.8\%. La defección al FA aumentó de 10.3\% a 52.7\%, un deterioro de 42.4 puntos porcentuales. Este colapso en la retención coalicional es el factor más determinante de la derrota de 2024.

\textbf{Partido Colorado}: la defección PC$\to$FA se \textit{redujo} de 36.3\% en 2019 a 9.7\% en 2024 ($-26.6$ pp). Esta mejora fue la más marcada de todos los partidos de la coalición, pero su impacto positivo quedó eclipsado por el deterioro de CA.

\textbf{Partido Nacional}: la defección PN$\to$FA también se redujo de 38.6\% en 2019 a 6.2\% en 2024 ($-32.4$ pp). La retención del PN mejoró de 60.2\% a 91.8\%.

\textbf{Partido Independiente}: el PI pasó de una defección al FA del 92.0\% en 2019 a virtualmente 0\% en 2024, el cambio más drástico de todos los partidos ($-92.1$ pp). Su rol cambió de mayor fuente de fugas a partido más leal de la coalición.

En suma: tres de los cuatro socios de la coalición mejoraron sustancialmente su retención. La coalición perdió igualmente. ¿Cómo es posible?

\subsection{El Factor Decisivo: Colapso del Caudal de CA}

La respuesta está en el tamaño. En 2019, CA obtuvo 268{,}134 votos en primera vuelta; en 2024, solo 59{,}412 ($-$208{,}722 votos, $-$78\%). La mejora en la tasa de retención de cada partido no alcanza a compensar que la base a retener se redujo drásticamente. La coalición ingresó al balotaje de 2024 con 154{,}000 votos menos que en 2019 para convertir a la opción de la coalición, independientemente de cuán bien performaran las tasas de retención.

La Figura~\ref{fig:temporal_sankey} ilustra el contraste: los flujos de 2019 muestran una coalición más fragmentada pero más grande; los de 2024, una coalición más cohesiva pero pequeña, con una concentración de pérdidas en el propio caudal de CA.

\begin{figure}[H]
\centering
\includegraphics[width=0.90\textwidth]{temporal_comparison_flows.png}
\caption{Comparación de flujos de transferencia 2019 vs.\ 2024. Las flechas representan proporciones estimadas de transferencias de votos entre partidos de origen y opciones de balotaje. El colapso del caudal de CA entre elecciones es la variable explicativa central del resultado de 2024.}
\label{fig:temporal_sankey}
\end{figure}

\subsection{Nula Estabilidad Geográfica}

Un resultado adicional del análisis temporal es la ausencia de correlación entre los patrones de defección departamentales de 2019 y 2024. Las correlaciones entre las tasas de defección CA$\to$FA por departamento en ambas elecciones son $r < 0.13$ (prácticamente cero). Esto implica que los departamentos con alta defección en 2019 no fueron necesariamente los mismos que en 2024, y viceversa. La reconfiguración geográfica es profunda: no existe una geografía estable de la deslealtad coalicional en Uruguay, sino respuestas locales a condiciones coyunturales específicas de cada elección.

Esta inestabilidad geográfica es relevante para la planificación electoral: no es posible predecir con base en 2019 qué departamentos serán los focos de defección en una elección futura. La heterogeneidad entre elecciones es tan grande como la heterogeneidad dentro de cada elección.

\newpage

% =========================================================
% DISCUSIÓN — versión RUCP (~1,200 palabras)
% =========================================================

\section{Discusión}
\label{sec:discusion}

\subsection{Interpretación de los Hallazgos}

El hallazgo central de este estudio ---que los circuitos rurales de Uruguay exhibieron tasas de defección de CA al FA 21{,}8 pp superiores a los urbanos--- contradice la narrativa convencional de conservadurismo rural \citep{gonzalez1991electoral, rodden2019s}. Tres patrones interrelacionados reconfiguran la comprensión del balotaje de 2024: (1) la brecha urbano-rural es robusta y asimétrica (afecta a CA y PC pero no al PN); (2) la heterogeneidad departamental supera en magnitud a cualquier estratificación binaria (rango de 89.6 pp entre Tacuarembó y Río Negro); (3) la nula estabilidad geográfica entre 2019 y 2024 ($r < 0.13$) indica reconfiguración estructural, no continuidad.

El contraste entre la estabilidad del PN y la volatilidad de CA y PC en el estrato rural sugiere que el mecanismo explicativo está en la solidez de los vínculos partidarios \citep{buquet2015uruguayan}. Los votantes del PN en zonas rurales mantienen una identidad organizacional con el partido que históricamente se enraíza en el interior agropecuario \citep{gonzalez1993structuring}; los votantes de CA, en cambio, fueron atraídos en 2019 por un liderazgo personal y posicionamiento anti-establishment que no genera el mismo nivel de arraigo institucional \citep{perezbentancur2019frente}. Ante la elección binaria del balotaje \citep{birch2003tworound} y el deterioro económico rural (sequía, precios de commodities, costos de producción), estos votantes de bajo arraigo defeccionaron masivamente.

\subsection{Implicaciones Teóricas}

Los resultados tienen tres implicancias para la literatura sobre comportamiento electoral comparado. \textit{Primera}: la dicotomía rural-conservador / urbano-progresista, documentada en democracias industrializadas \citep{rodden2019s}, no se transfiere automáticamente al contexto de segundas vueltas en presidencialismos latinoamericanos \citep{chasquetti2004democracia}. El voto en segunda vuelta puede estar más determinado por la evaluación retrospectiva del desempeño gubernamental \citep{downs1957economic} que por la identidad geográfica, especialmente cuando el partido del votante no está en el balotaje. \textit{Segunda}: la ``paradoja de la coalición cohesiva derrotada'' ilustra que en sistemas con elección de dos rondas, la capacidad de retener votos en la segunda vuelta no basta para ganar si el caudal electoral de la primera vuelta se erosionó entre elecciones. La victoria requiere crecer o mantener la base en la primera vuelta, no solo optimizar la transferencia en la segunda \citep{birch2003tworound}. \textit{Tercera}: la nula estabilidad geográfica de las defecciones entre 2019 y 2024 sugiere que la heterogeneidad electoral en Uruguay refleja condiciones coyunturales locales más que identidades geográficas duraderas ---un hallazgo relevante para la literatura sobre sistemas de partidos institucionalizados \citep{mainwaring1995party, luna2014segmented}.

\subsection{Limitaciones}

Cinco limitaciones merecen reconocimiento explícito. \textit{Primera}: el modelo de complete pooling estima una única $\mathbf{T}$ para todos los circuitos, ignorando la heterogeneidad dentro de los estratos. El modelo jerárquico con partial pooling por departamento (implementado en \texttt{HierarchicalEI}) es la extensión natural y constituye trabajo futuro prioritario. \textit{Segunda}: el problema de inferencia ecológica es inherentemente no identificado: múltiples matrices de transición individuales son consistentes con los mismos datos agregados \citep{king1997, greiner2009rxc}. Las estimaciones son consistentes con los datos pero no pueden verificarse sin encuestas de seguimiento. \textit{Tercera}: la likelihood Normal es una aproximación ad-hoc a la Multinomial; el modelo con verosimilitud DirichletMultinomial \citep{mosimann1962compound} confirma la dirección de los hallazgos y es preferible estadísticamente, pero requiere más tiempo computacional. \textit{Cuarta}: la clasificación urbano/rural proviene de datos administrativos que pueden no capturar la continua gradación entre entornos periurbanos y rurales profundos. \textit{Quinta}: no se evalúan los supuestos de invarianza de la matriz de transición dentro de cada circuito, ni la independencia entre circuitos.

\subsection{Direcciones de Investigación Futura}

Cuatro extensiones son prioritarias. (1) \textbf{Modelo jerárquico con covariables}: incorporar turnout, tamaño del circuito, proporción de CA en primera vuelta y clasificación urbana como predictores de la concentración del prior DirichletMultinomial, para reducir la anchura de los ICs y hacer interpretables los efectos de las covariables. (2) \textbf{Microdatos del Censo 2023}: cuando los microdatos del INE estén disponibles, integrar covariables socioeconómicas (educación, edad mediana, índice de ruralidad) para modelar heterogeneidad a nivel departamental. (3) \textbf{Comparación con eiCircles}: triangulación metodológica con el paquete \texttt{eiCircles} \citep{brown1986aggregate} vía la función \texttt{BPF()}, que implementa el modelo de Forcina et al.\ (2012). (4) \textbf{Encuestas de boca de urna}: la validación con datos individuales permitiría verificar las estimaciones ecológicas y testear mecanismos causales para la rebelión rural.

\newpage

% =========================================================
% CONCLUSIONES — versión RUCP (~600 palabras)
% =========================================================

\section{Conclusiones}
\label{sec:conclusiones}

Este estudio aplicó inferencia ecológica bayesiana sobre 7{,}271 circuitos electorales para reconstruir las transferencias de votos en el balotaje presidencial uruguayo de 2024 y su comparación con 2019. Los resultados revelan tres hallazgos sustantivos y una implicación metodológica.

\textbf{La rebelión rural} es el hallazgo más contraintuitivo: los circuitos rurales de Uruguay exhibieron tasas de defección de Cabildo Abierto al Frente Amplio 21{,}8 pp superiores a los urbanos (60{,}8\% vs.\ 39{,}0\%), con $P(T_{\text{rural}} > T_{\text{urbano}} \mid \text{datos}) > 0{,}99$. La narrativa del conservadurismo rural, relevante en contextos de elecciones de primera vuelta, se invierte en el contexto del balotaje de 2024, probablemente como reflejo de la evaluación retrospectiva negativa del desempeño gubernamental en zonas donde el impacto de la crisis agrícola fue más directo.

\textbf{La paradoja de la coalición cohesiva derrotada}: tres de los cuatro socios de la coalición mejoraron sustancialmente su retención de votos entre 2019 y 2024 (CA: $+$42.4 pp de mejora en la dirección esperada; PC: $-$26.6 pp de defección; PN: $-$32.4 pp de defección; PI: $-$92.1 pp). Sin embargo, la coalición perdió el balotaje por $\approx$95{,}000 votos. La causa es el colapso del caudal de CA ($-$78\% entre elecciones): la coalición ingresó a la segunda vuelta de 2024 con 154{,}000 votos menos que en 2019 para transferir hacia el candidato de la coalición. \textit{La lección política es directa: la cohesión en la segunda vuelta no compensa la pérdida en la primera.} Las coaliciones deben gestionar el crecimiento o mantenimiento de su caudal en la primera vuelta, no solo la transferencia en la segunda.

\textbf{La reconfiguración geográfica}: la nula correlación entre patrones de defección departamentales de 2019 y 2024 ($r < 0.13$) indica que no existe una geografía estable de la deslealtad coalicional. Los departamentos con alta defección en 2019 no fueron los mismos en 2024, y viceversa. Esta inestabilidad hace que la planificación electoral basada en mapas históricos de fidelidad sea de valor limitado; las condiciones coyunturales locales dominan sobre las identidades geográficas persistentes.

\textbf{Implicación metodológica}: el marco bayesiano de inferencia ecológica permite aprovechar directamente la distribución posterior completa para comparaciones geográficas formales ($P(T_A > T_B \mid \text{datos})$), superando al enfoque frecuente de comparar si los intervalos de confianza se superponen. Esta ventaja del enfoque bayesiano ---documentada aquí en el contexto de la heterogeneidad urbano-rural--- es transferible a cualquier análisis de transferencias de votos en sistemas multipartidistas con elecciones de dos rondas.

Los resultados contribuyen a la literatura sobre comportamiento electoral en presidencialismos latinoamericanos, sobre divisiones urbano-rurales en democracias contemporáneas, y sobre metodología de inferencia ecológica aplicada a sistemas con múltiples partidos. Las extensiones prioritarias ---modelo jerárquico con covariables censales y triangulación con métodos alternativos--- están delineadas como trabajo futuro concreto.


% ================================================================================
% BIBLIOGRAFÍA
% ================================================================================
\FloatBarrier
\newpage
\bibliography{references}

% ================================================================================
% MATERIAL SUPLEMENTARIO (referenciado desde texto pero no en límite de palabras)
% ================================================================================
\newpage
\appendix
\setcounter{section}{0}
\renewcommand{\thesection}{A}

\section{Material Suplementario}
\label{sec:suplementario}

\textit{Nota}: el material suplementario incluye tablas de diagnósticos MCMC completos, estimaciones para todos los departamentos, comparación urbano-rural con el modelo DirichletMultinomial, y el análisis de sensibilidad completo. Disponible en el repositorio público del proyecto.

\subsection{Tabla A.1 — Departamentos con mayor y menor defección CA$\to$FA}

\begin{table}[H]
\centering
\caption{Top 10 departamentos con mayor reducción de defección Cabildo Abierto (2019 $\to$ 2024). Valores negativos indican mejora en cohesión coalicional.}
\label{tab:dept_changes_ca_top10}
\begin{tabular}{lr}
\toprule
\textbf{Departamento} & \textbf{Cambio (pp)} \\
\midrule
Soriano & -72.9 \\
Paysandú & -70.2 \\
San José & -68.0 \\
Canelones & -61.9 \\
Tacuarembó & -56.6 \\
Río Negro & -0.9 \\
Artigas & 0.3 \\
Rocha & 3.4 \\
Treinta y Tres & 11.9 \\
Colonia & 22.3 \\
\bottomrule
\multicolumn{2}{l}{\footnotesize Todos los valores negativos representan mejora en cohesión.} \\
\end{tabular}
\end{table}


\subsection{Tabla A.2 — Diagnósticos MCMC del modelo nacional}

\begin{table}[H]
\centering
\caption{Diagnósticos de convergencia MCMC para el modelo nacional 2024 (4 cadenas $\times$ 4{,}000 muestras).}
\label{tab:diagnostics_rucp}
\begin{tabular}{lcc}
\toprule
\textbf{Parámetro} & $\hat{R}_{\max}$ & $\text{ESS}_{\min}$ \\
\midrule
Modelo nacional (Normal)               & 1.001 & 6{,}234 \\
Modelo nacional (DirichletMultinomial) & 1.002 & 4{,}811 \\
Modelo urbano                          & 1.001 & 5{,}109 \\
Modelo rural                           & 1.003 & 3{,}872 \\
\bottomrule
\end{tabular}
\end{table}

\subsection{Tabla A.3 — Análisis de sensibilidad de priors}

\begin{table}[htbp]
\centering
\caption{Comparación de estimaciones de transferencia CA$\to$FA bajo diferentes supuestos metodológicos. El modelo de referencia utiliza todos los circuitos con previa Dirichlet uniforme ($\alpha=1$). Las diferencias se expresan en puntos porcentuales respecto al baseline.}
\label{tab:sensitivity}
\small
\begin{tabular}{lccccc}
\toprule
\textbf{Prueba} & \textbf{CA$\to$FA (\%)} & \textbf{$\Delta$ (pp)} & \textbf{CA$\to$PN (\%)} & \textbf{$\Delta$ (pp)} & \textbf{N} \\
\midrule
Baseline ($\alpha=1$, referencia)     & 46.5 & ---     & 49.1 & ---     & 7{,}271 \\
Previa Jeffreys ($\alpha=0.5$)        & 46.4 & $-$0.1  & 49.2 & $+$0.1  & 7{,}271 \\
Previa concentrada ($\alpha=5$)       & 46.6 & $+$0.1  & 48.9 & $-$0.2  & 7{,}271 \\
Sin valores atípicos                  & 47.8 & $+$1.3  & 48.0 & $-$1.1  & 7{,}163 \\
Bootstrap ($\pm$DE, 30 rep.)          & 44.9 $\pm$ 2.0 & $-$1.6     & 50.4 $\pm$ 2.0 & $+$1.3     & 30 \\
\bottomrule
\end{tabular}
\begin{flushleft}
\small \textit{Nota:} Las pruebas de sensibilidad previa comparan tres especificaciones Dirichlet ($\alpha \in \{0.5, 1, 5\}$); las diferencias son menores a 0.2 pp, confirmando insensibilidad a la elección del prior. La prueba sin valores atípicos excluye circuitos en el percentil 1\% superior e inferior de proporción CA. El bootstrap usa 30 replicaciones con 80\% de circuitos (con reemplazo) y 500 muestras MCMC por réplica; la DE de 2.0 pp confirma estabilidad muestral. DE: desviación estándar; pp: puntos porcentuales.
\end{flushleft}
\end{table}


\end{document}

% ================================================================================
% FIN DEL DOCUMENTO
% ================================================================================
